\section{Related Work}
\label{sec:related}

\subsection{Cooperative Perception}

Cooperative perception architectures trade off bandwidth against
perceptual fidelity.  Early fusion (raw data exchange) is
bandwidth-prohibitive.  Late fusion (object-level sharing) discards
spatial context needed for occluded-object detection.  Intermediate
fusion dominates recent work: Where2Comm~\cite{hu2022where2comm}
uses spatial-confidence attention maps,
V2X-ViT~\cite{v2xvit2022} applies vision transformers to BEV
features, and CoBEVT~\cite{xu2022cobevt} fuses multi-scale BEV
representations with axial attention.

CMP~\cite{cmp2024,cmp_ieee_2023} is the only prior system that
unifies cooperative perception, multi-object tracking (AB3DMOT), and
trajectory prediction (MTR) into a single evaluated pipeline.  CMP
reports results at 2--7 agents on OPV2V and V2V4Real, with no
deadline constraint.  Our work extends the same pipeline family to
4--64 agents, characterizes the resulting deadline cliff, and
introduces adaptation mechanisms that maintain feasibility at dense
scale.

\subsection{Edge-Hosted Perception Systems}

EMP~\cite{EMP} demonstrated edge-hosted LiDAR fusion at 93\,ms
end-to-end on a real intersection, but evaluates detection only (no
tracking or prediction).  Harbor~\cite{harbor} tests a hybrid
V2I/V2V architecture on a 3-vehicle Mcity deployment, using frame
timeouts rather than a compute-budget controller.
F-Cooper~\cite{fcooper} evaluates point-cloud feature sharing
at the edge but does not address the computational scaling problem.
EdgeCooper~\cite{edgecooper2023} considers communication-aware
scheduling for edge cooperative perception but evaluates at
$N \leq 5$.

None of these systems evaluate the full pipeline (fusion + tracking +
prediction) at $N > 10$, nor do they model the joint-pipeline deadline
cliff that emerges from cross-stage compute coupling.

\subsection{Adaptive Compute for Autonomous Vehicles}

Single-vehicle adaptive compute has been studied in the context of
streaming perception~\cite{li2020streaming} and modular
self-driving platforms~\cite{gog2022pylot}.  These approaches
adjust model complexity or scheduling policy per frame based on
deadline proximity.  They operate within a single vehicle's compute
budget and do not address the multi-agent cooperative setting where
$N$ agents share one edge deadline.

\subsection{Multi-Object Tracking}

Classical multi-object tracking uses Kalman filtering with Hungarian
association (SORT~\cite{Bewley2016_sort},
AB3DMOT~\cite{ab3dmotTracker2020}).  Recent work explores learned
association and state-space models for temporal
dynamics~\cite{mambatrack2024, trackssm2024}.
We use AB3DMOT with vectorized Kalman operations (1.8$\times$ speedup
over the original per-track loop), which provides sufficient tracking
quality while remaining the cheapest pipeline stage.

\subsection{Trajectory Prediction}

MTR~\cite{shi2022mtr} and MTR++~\cite{shi2023mtrpp} use transformer
decoders with motion-query attention to produce multi-modal trajectory
forecasts.
SMART~\cite{wu2024smart} tokenizes trajectories for autoregressive
generation.  We evaluated both on OPV2V data: MTR achieves 4.5\,m ADE
at 57\,ms, while SMART yields 42.7\,m ADE at 137\,ms (reflecting a
Waymo-to-CARLA domain gap and exceeding the 100\,ms deadline on its
own).  MTR is the only viable option for deadline-constrained edge
inference in our pipeline.
